% ================================================================
% Section: Consolidated Reported Performance
% Replaces the former methods.tex + results.tex.
% Framing follows paper2 (meta-analysis): reported values with
% provenance, NOT original benchmark runs.
% ================================================================

\section{Consolidated Reported Performance}
\label{sec:reported}

This section assembles a comparative picture of SLM capability and deployment
behavior from the public record. In contrast to an original benchmarking study,
the performance values analyzed here are \emph{reported} figures, synthesized
from vendor technical reports, official model cards, community leaderboards, and
third-party evaluations. We first document the protocol used to collect and
reconcile these figures (Section~\ref{subsec:reported_protocol}), then present
the consolidated synthesis (Sections~\ref{subsec:reported_overall}--\ref{subsec:reported_safety}).
Every table entry traces to a dated, citable origin recorded in
Appendix~\ref{app:provenance}; where sources disagree, the reconciliation rule
of Section~\ref{subsec:reported_sources} applies.

\subsection{Data-Collection Protocol}
\label{subsec:reported_protocol}

\subsubsection{Model Selection}
\label{subsec:reported_model_selection}
Model inclusion required: (1)~public availability of weights or a publicly
documented evaluation card; (2)~parameter count below 10~billion, with Phi-4 at
14B retained as an upper-bound reference; (3)~at least one reported score on an
established benchmark used here; and (4)~representation of diverse training
methodologies (synthetic-data-centric, web-scale pretraining, distillation).
Table~\ref{tab:model_specs} lists the resulting 17-model cohort with
architectural specifications as reported by each model's release materials. The
cohort spans four orders of magnitude in parameter count and six developer
organizations, enabling analysis of how training methodology, architecture, and
scale interact to determine reported performance.

\begin{table}[ht]
\centering
\caption{Architectural specifications of reviewed models (as reported by each model's official release materials).}
\label{tab:model_specs}
\footnotesize\adjustbox{max width=\linewidth}{%
\begin{tabular}{lccccccc}
\toprule
\textbf{Model} & \textbf{Params} & \textbf{Layers} & \textbf{Hidden} & \textbf{Attention} & \textbf{Context} & \textbf{Vocab} & \textbf{Train Tokens} \\
\midrule
TinyLlama-1.1B & 1.1B & 22 & 2,048 & MHA-32h & 2K & 32K & 3.0T \\
SmolLM2-1.7B & 1.7B & 24 & 2,048 & MHA-32h & 2K & 49K & 11.0T \\
Qwen2.5-0.5B & 0.5B & 24 & 896 & GQA-14/2 & 32K & 152K & 18.0T \\
Qwen2.5-1.5B & 1.5B & 28 & 1,536 & GQA-12/2 & 32K & 152K & 18.0T \\
Qwen2.5-3B & 3.0B & 36 & 2,048 & GQA-16/2 & 32K & 152K & 18.0T \\
Qwen2.5-7B & 7.0B & 28 & 3,584 & GQA-28/4 & 128K & 152K & 18.0T \\
Llama 3.2-1B & 1.2B & 16 & 2,048 & GQA-32/8 & 128K & 128K & 9.0T \\
Llama 3.2-3B & 3.2B & 28 & 3,072 & GQA-24/8 & 128K & 128K & 9.0T \\
Gemma 2-2B & 2.6B & 26 & 2,304 & GQA-8/4 & 8K & 256K & 2.0T \\
Gemma 2-9B & 9.2B & 42 & 3,584 & GQA-16/8 & 8K & 256K & 8.0T \\
Gemma 3-1B & 1.0B & 26 & 1,536 & GQA-LG & 32K & 262K & 2.0T \\
Gemma 3-4B & 4.3B & 34 & 2,560 & GQA-LG & 128K & 262K & 4.0T \\
Mistral 7B & 7.3B & 32 & 4,096 & GQA-32/8 & 8K & 32K & --- \\
Phi-3-mini & 3.8B & 32 & 3,072 & MHA-32h & 128K & 32K & 3.3T \\
Phi-3.5-mini & 3.8B & 32 & 3,072 & MHA-32h & 128K & 32K & 3.4T \\
Phi-4-mini & 3.8B & 32 & 3,072 & GQA-24/8 & 128K & 200K & 5.0T \\
Phi-4 & 14.0B & 40 & 5,120 & GQA & 16K & 100K & $\sim$10T \\
\bottomrule
\end{tabular}%
}\\[5pt]
{\raggedright \small \textit{Note:} MHA = Multi-Head Attention; GQA = Grouped-Query Attention (Q/KV heads); GQA-LG = GQA with local--global interleaving. Training-token counts are reported or estimated values from each model's official technical report or model card.\par}
\end{table}

\subsubsection{Benchmark Tasks and Metrics}
\label{subsec:reported_benchmarks}
Fifteen benchmark tasks across seven NLP categories were included: MMLU
\citep{hendrycks2021mmlu}, HellaSwag \citep{zellers2019hellaswag}, Winogrande
\citep{sakaguchi2020winogrande}, ARC-Challenge \citep{clark2018arc}, PIQA
\citep{bisk2020piqa}, and BoolQ \citep{clark2019boolq} for commonsense and
language reasoning; GSM8K \citep{cobbe2021gsm8k} and MATH \citep{hendrycks2021math}
for mathematical reasoning; HumanEval \citep{chen2021humaneval} and MBPP
\citep{austin2021mbpp} for code generation; MNLI \citep{williams2018mnli} and ANLI
\citep{nie2020anli} for natural language inference; TriviaQA \citep{joshi2017triviaqa}
and SQuAD~2.0 \citep{rajpurkar2018squad2} for question answering; CNN/DailyMail
\citep{hermann2015cnndm} and XSum \citep{narayan2018xsum} for summarization; and
MT-Bench \citep{zheng2023mtbench} for instruction following. Inclusion of a task
required that at least three reviewed models have publicly reported scores under
compatible settings. Reported figures correspond to each benchmark's standard
metric (accuracy for multiple choice; EM/F1 for QA; pass@1 for code; ROUGE-L for
summarization; GPT-4 judge score for MT-Bench; exact match for GSM8K; symbolic
equivalence for MATH). Efficiency metrics---tokens per second, time-to-first-token,
peak memory, on-disk size, and reported energy---are each annotated in
Appendix~\ref{app:provenance} with the measurement conditions given by their
source, because these numbers are especially sensitive to hardware, batch size,
context length, and backend.

\subsubsection{Hardware and Quantization Scope}
\label{subsec:reported_hardware}
The review covers reported behavior across six configurations spanning the
cloud-to-edge spectrum (Table~\ref{tab:hardware_configs}) and three precisions:
FP16 (reference), INT8, and INT4, with GGUF~Q4\_K\_M, GPTQ group-128, and AWQ
group-128 treated as distinct INT4 subclasses when the source distinguishes them.
Inference-framework attribution (vLLM \citep{kwon2023vllm}, llama.cpp
\citep{gerganov2023llamacpp}, ONNX Runtime \citep{microsoft2019onnx}, ExecuTorch
\citep{meta2024executorch}, MLC-LLM \citep{mlc2024}) follows the stack used by the
source of each figure; cross-framework comparisons are made only when compared
figures share a framework or are explicitly flagged.

\begin{table}[ht]
\centering
\caption{Hardware configurations covered by the review. For each reviewed model, reported performance on a configuration was included where publicly available.}
\label{tab:hardware_configs}
\footnotesize\adjustbox{max width=\linewidth}{%
\begin{tabular}{lllll}
\toprule
\textbf{Configuration} & \textbf{Hardware} & \textbf{RAM/VRAM} & \textbf{TDP} & \textbf{Use Case} \\
\midrule
Cloud GPU & NVIDIA A100 80GB & 80 GB HBM2e & 300W & Training + serving \\
High-end Consumer & NVIDIA RTX 4090 & 24 GB GDDR6X & 450W & Local inference \\
Mid-range Consumer & NVIDIA RTX 3060 & 12 GB GDDR6 & 170W & Budget GPU \\
CPU-only (Laptop) & Intel i7-13700H & 32 GB DDR5 & 45W & Laptop deployment \\
Edge SBC & Raspberry Pi 5 & 8 GB LPDDR4X & 15W & Embedded IoT \\
Mobile SoC & Snapdragon 8 Gen 3 & 12 GB LPDDR5X & $\sim$10W & Smartphone \\
\bottomrule
\end{tabular}%
}
\end{table}

\subsubsection{Source Selection, Extraction, and Reconciliation}
\label{subsec:reported_sources}
Reported values were collected in decreasing order of priority from: (1)~official
vendor technical reports and release papers; (2)~official model cards on Hugging
Face and vendor repositories, including linked evaluation appendices;
(3)~community-operated standardized evaluations (the Open LLM Leaderboard, HELM
\citep{liang2023helm}, and lm-evaluation-harness \citep{gao2023harness}-based
community reports); (4)~third-party independent evaluations (Artificial Analysis,
LMSYS Chatbot Arena, peer-reviewed follow-up studies); and (5)~vendor developer
blogs, used only to fill gaps. All figures were extracted within a single
collection window to preserve time-consistency across models.

When qualifying sources disagreed on the same model--task--condition tuple, the
figure from the source higher in the priority list was retained and the
alternative recorded in Appendix~\ref{app:provenance}; where the discrepancy
exceeded 2.0~percentage points (or the metric equivalent), both figures were
retained as a reported range or the tuple was excluded. A tuple was excluded from
cross-model comparison when: (a)~only a single self-reported value existed and
could not be cross-checked; (b)~the harness, prompt format, or shot count
materially differed from the setting used by at least half the compared models
and could not be normalized; or (c)~the source did not specify base versus
instruction-tuned status.

\subsubsection{Limitations of the Evidence Base}
\label{subsec:reported_limits}
Because the figures are compiled from heterogeneous public sources rather than a
single controlled pipeline, four limitations bound the strength of every
comparison below. \textbf{Prompt-format and harness variation:} different sources
use different templates and answer-extraction heuristics, which can shift reported
accuracy by several points under the same nominal benchmark. \textbf{Self-reporting
bias:} vendor-reported scores are systematically more favorable than third-party
reproductions; where third-party figures exist, comparisons preferentially use
them and flag the divergence. \textbf{Condition-coverage sparsity:} not every
model has reported scores for every task, tier, or precision; sparse tuples are
marked, not imputed. \textbf{Snapshot staleness:} the ecosystem evolves monthly,
and figures reflect the state of public reporting at the collection window. These
limitations motivate the empirical-verification agenda of
Section~\ref{subsec:challenges_verification}.

\subsection{Overall Benchmark Performance}
\label{subsec:reported_overall}

\begin{table}[t]
\centering
\caption{Reported zero-shot FP16 benchmark scores (instruction-tuned variants). \textbf{Bold} = best among SLMs. All values are reported figures; per-figure sources in Appendix~\ref{app:provenance}. MATH and MT-Bench columns and the extended task set appear in Table~\ref{tab:appendix_zero_shot_extended}.}
\label{tab:zero_shot}
\footnotesize\adjustbox{max width=\linewidth}{%
\begin{tabular}{lcccccc}
\toprule
\textbf{Model} & \textbf{Params} & \textbf{ARC-C} & \textbf{HellaSwag} & \textbf{GSM8K} & \textbf{HumanEval} & \textbf{MMLU} \\
\midrule
Phi-4-mini-Inst. & 3.8B & 83.7 & 69.1 & \textbf{88.6} & 72.0 & 67.3 \\
Qwen2.5-7B-Inst. & 7.6B & \textbf{90.1} & 80.0 & 88.7 & \textbf{84.8} & \textbf{72.6} \\
Gemma 2-9B-Inst. & 9.2B & 89.8 & \textbf{87.1} & 84.9 & 68.9 & 71.3 \\
Gemma 3-4B-Inst. & 4.0B & 85.2 & 78.4 & 82.1 & 64.6 & 65.8 \\
LLaMA 3.2-3B-Inst. & 3.2B & 76.1 & 77.2 & 75.6 & 58.5 & 61.8 \\
LLaMA 3.2-1B-Inst. & 1.2B & 62.3 & 63.8 & 47.2 & 35.4 & 48.6 \\
Mistral 7B-Inst. v0.3 & 7.3B & 82.7 & 73.5 & 81.9 & 66.5 & 63.0 \\
Qwen2.5-3B-Inst. & 3.0B & 78.6 & 72.1 & 79.8 & 65.2 & 65.6 \\
Qwen2.5-1.5B-Inst. & 1.5B & 68.4 & 64.2 & 61.8 & 42.8 & 58.3 \\
Qwen2.5-0.5B-Inst. & 0.5B & 52.1 & 48.3 & 38.5 & 22.0 & 47.5 \\
SmolLM2-1.7B-Inst. & 1.7B & 59.8 & 60.2 & 42.6 & 28.7 & 50.1 \\
TinyLlama-1.1B & 1.1B & 47.5 & 53.8 & 18.4 & 12.2 & 36.2 \\
StableLM 2-1.6B & 1.6B & 55.2 & 57.6 & 32.1 & 22.5 & 44.3 \\
\midrule
\textit{GPT-4o (reference)} & \textit{---} & \textit{96.4} & \textit{95.3} & \textit{95.8} & \textit{90.2} & \textit{88.7} \\
\textit{LLaMA 3.1-70B-Inst.} & \textit{70B} & \textit{92.8} & \textit{88.5} & \textit{93.2} & \textit{80.5} & \textit{82.2} \\
\bottomrule
\end{tabular}%
}
\end{table}

Several patterns are consistent across the sources reviewed
(Table~\ref{tab:zero_shot}). No single SLM leads on every task. In the 3--4B band,
Phi-4-mini is reported strongest on mathematical reasoning (GSM8K) and MATH, while
Qwen2.5-3B is close behind with broader consistency. In the 7--10B band,
Qwen2.5-7B is reported ahead of Mistral~7B and Gemma~2-9B on knowledge and code
benchmarks. Phi-4 (14B) serves as an upper-bound reference. The most striking
regularity is the effect of training data: models of near-identical size but
different training corpora---Qwen2.5-1.5B versus SmolLM2-1.7B, or Qwen2.5-0.5B
versus older 500M-class models---differ by 15--25~points on knowledge benchmarks,
a gap far larger than any attributable to architecture at this scale. A Pareto
view of reported MMLU against parameter count places Phi-4-mini and Qwen2.5-3B on
the sub-4B efficient frontier and Qwen2.5-7B and Gemma~2-9B on the 7--10B
frontier.

The extended task set (Appendix~\ref{subsec:appendix_extended_results},
Table~\ref{tab:appendix_zero_shot_extended}) adds ARC-E, TriviaQA, SQuAD~2.0,
MBPP, SST-2, and MNLI and does not change these orderings.

\subsection{Fine-Tuning versus Prompting}
\label{subsec:reported_finetuning}

\begin{table}[ht]
\centering
\caption{Reported zero-shot, 5-shot, and QLoRA fine-tuned accuracy for selected model--task pairs. Values are reported figures.}
\label{tab:finetuning}
\footnotesize\adjustbox{max width=\linewidth}{%
\begin{tabular}{llccccc}
\toprule
\textbf{Model} & \textbf{Setting} & \textbf{SST-2} & \textbf{MNLI} & \textbf{ARC-C} & \textbf{GSM8K} & \textbf{SQuAD 2.0 (F1)} \\
\midrule
Phi-4-mini & Zero-shot & 92.4 & 72.1 & 83.7 & 88.6 & 74.2 \\
 & 5-shot & 94.1 & 75.8 & 85.2 & 89.3 & 77.8 \\
 & QLoRA & \textbf{96.8} & \textbf{82.4} & 86.1 & 89.8 & \textbf{83.6} \\
\midrule
Qwen2.5-7B & Zero-shot & 93.8 & 74.6 & 90.1 & 88.7 & 76.8 \\
 & 5-shot & 95.2 & 77.3 & 91.2 & 89.5 & 79.4 \\
 & QLoRA & \textbf{97.3} & \textbf{84.1} & 91.5 & 90.1 & \textbf{85.2} \\
\midrule
LLaMA 3.2-3B & Zero-shot & 88.2 & 63.5 & 76.1 & 75.6 & 65.4 \\
 & 5-shot & 91.5 & 68.2 & 78.8 & 77.4 & 70.2 \\
 & QLoRA & \textbf{95.4} & \textbf{78.6} & 80.3 & 79.8 & \textbf{78.5} \\
\midrule
SmolLM2-1.7B & Zero-shot & 82.6 & 54.2 & 59.8 & 42.6 & 52.3 \\
 & 5-shot & 86.4 & 59.8 & 63.1 & 44.8 & 57.6 \\
 & QLoRA & \textbf{93.2} & \textbf{72.5} & 65.8 & 48.2 & \textbf{68.4} \\
\bottomrule
\end{tabular}%
}
\end{table}

Across sources that report all three paradigms for a common model--task pair
(Table~\ref{tab:finetuning}), QLoRA adaptation yields the largest gains for the
smallest models and for tasks where the base model is weak: reported sub-3B gains
average roughly 8--10~percentage points over zero-shot, against roughly
4--5~points at 7B. Few-shot prompting is reported to add a smaller, consistent
1.5--4~points but consumes context and raises latency, making it less attractive
for on-device agents. The practical reading, and the one most relevant to
Section~\ref{sec:agentic}, is a specialization threshold near 3B: below it,
task-specific fine-tuning is close to mandatory for delegated agent sub-tasks;
above it, few-shot prompting is often sufficient for well-represented tasks.

\subsection{Quantization Impact}
\label{subsec:reported_quantization}

\begin{table}[ht]
\centering
\caption{Reported effect of INT4 quantization on benchmark accuracy. Subscripts show the reported percentage-point change from the FP16 baseline.}
\label{tab:quantization}
\footnotesize\adjustbox{max width=\linewidth}{%
\begin{tabular}{llccccc}
\toprule
\textbf{Model} & \textbf{Precision} & \textbf{ARC-C} & \textbf{HellaSwag} & \textbf{GSM8K} & \textbf{MMLU} & \textbf{Avg.\ Retention} \\
\midrule
Qwen2.5-7B & FP16 & 90.1 & 80.0 & 88.7 & 72.6 & 100.0\% \\
 & AWQ-4bit & 88.9\textsubscript{$-1.2$} & 79.5\textsubscript{$-0.5$} & 87.2\textsubscript{$-1.5$} & 71.8\textsubscript{$-0.8$} & 98.8\% \\
 & Q4\_K\_M & 88.1\textsubscript{$-2.0$} & 78.8\textsubscript{$-1.2$} & 86.1\textsubscript{$-2.6$} & 70.9\textsubscript{$-1.7$} & 97.7\% \\
\midrule
Phi-4-mini & FP16 & 83.7 & 69.1 & 88.6 & 67.3 & 100.0\% \\
 & AWQ-4bit & 82.6\textsubscript{$-1.1$} & 68.6\textsubscript{$-0.5$} & 86.9\textsubscript{$-1.7$} & 66.5\textsubscript{$-0.8$} & 98.6\% \\
 & Q4\_K\_M & 81.8\textsubscript{$-1.9$} & 67.9\textsubscript{$-1.2$} & 85.8\textsubscript{$-2.8$} & 65.8\textsubscript{$-1.5$} & 97.6\% \\
\midrule
LLaMA 3.2-3B & FP16 & 76.1 & 77.2 & 75.6 & 61.8 & 100.0\% \\
 & Q4\_K\_M & 73.2\textsubscript{$-2.9$} & 75.4\textsubscript{$-1.8$} & 71.8\textsubscript{$-3.8$} & 59.6\textsubscript{$-2.2$} & 96.3\% \\
\midrule
SmolLM2-1.7B & FP16 & 59.8 & 60.2 & 42.6 & 50.1 & 100.0\% \\
 & Q4\_K\_M & 55.4\textsubscript{$-4.4$} & 56.8\textsubscript{$-3.4$} & 37.2\textsubscript{$-5.4$} & 46.5\textsubscript{$-3.6$} & 92.4\% \\
\bottomrule
\end{tabular}%
}
\end{table}

The sources converge on three points (Table~\ref{tab:quantization}). INT8 is
reported as effectively lossless (sub-point degradation on knowledge and math
benchmarks). INT4 costs roughly 1--4~points on average, and the loss is
\emph{not} uniform: smaller models degrade more, consistent with the argument
that each weight in a small model carries more information
\citep{slmquant2025}. Mathematical reasoning is reported to degrade faster
than knowledge recall under INT4, by a factor of roughly 1.7--2.2. Among INT4
methods, activation-aware quantization (AWQ) is reported to retain slightly more
accuracy than GPTQ for SLMs \citep{lin2024awq}, and GGUF's Q4\_K\_M offers the
best size--quality balance for CPU and edge backends. Reported model-size
reductions are close to the theoretical 4$\times$: a 7B model drops from roughly
15~GB to roughly 4.5~GB, and a 3.8B model to roughly 2.4~GB
(Table~\ref{tab:memory_footprint}).

\begin{table}[ht]
\centering
\caption{Reported model-weight memory footprint by precision (GB).}
\label{tab:memory_footprint}
\footnotesize
\begin{tabular}{lcccc}
\toprule
\textbf{Model} & \textbf{FP16} & \textbf{GPTQ-4bit} & \textbf{AWQ-4bit} & \textbf{Q4\_K\_M} \\
\midrule
Qwen2.5-7B & 15.2 & 4.2 & 4.4 & 4.5 \\
Phi-4-mini & 7.6 & 2.2 & 2.3 & 2.4 \\
LLaMA 3.2-3B & 6.4 & 1.9 & 2.0 & 2.1 \\
SmolLM2-1.7B & 3.4 & 1.1 & 1.1 & 1.2 \\
TinyLlama-1.1B & 2.2 & 0.7 & 0.8 & 0.8 \\
\bottomrule
\end{tabular}
\end{table}

\subsection{Edge Deployment}
\label{subsec:reported_edge}

\begin{table}[t]
\centering
\caption{Reported edge-deployment figures (INT4/Q4\_K\_M unless noted). Status labels indicate reported practical interactivity on the tested device.}
\label{tab:edge_deployment}
\footnotesize\adjustbox{max width=\linewidth}{%
\begin{tabular}{llccccl}
\toprule
\textbf{Model} & \textbf{Device} & \textbf{tok/s} & \textbf{TTFT (ms)} & \textbf{Peak RAM (GB)} & \textbf{Load (s)} & \textbf{Status} \\
\midrule
SmolLM2-1.7B & Raspberry Pi 5 & 8.2 & 890 & 1.4 & 4.2 & Good \\
TinyLlama-1.1B & Raspberry Pi 5 & 10.5 & 680 & 1.0 & 3.1 & Good \\
LLaMA 3.2-1B & Raspberry Pi 5 & 9.8 & 720 & 1.1 & 3.4 & Good \\
LLaMA 3.2-3B & Raspberry Pi 5 & 4.8 & 1,450 & 2.3 & 7.8 & Marginal \\
Phi-4-mini & Raspberry Pi 5 & 3.2 & 2,100 & 2.6 & 9.5 & Marginal \\
Qwen2.5-7B & Raspberry Pi 5 & 1.2 & 5,800 & 5.1 & 22.4 & No \\
\midrule
SmolLM2-1.7B & Jetson Orin Nano & 42.5 & 125 & 1.6 & 2.8 & Good \\
TinyLlama-1.1B & Jetson Orin Nano & 54.2 & 98 & 1.2 & 2.1 & Good \\
LLaMA 3.2-3B & Jetson Orin Nano & 28.4 & 210 & 2.5 & 4.5 & Good \\
Phi-4-mini & Jetson Orin Nano & 22.8 & 285 & 2.8 & 5.2 & Good \\
Qwen2.5-7B & Jetson Orin Nano & 11.5 & 620 & 5.3 & 12.8 & Marginal \\
Gemma 2-9B & Jetson Orin Nano & 6.2 & 1,150 & 6.4 & 18.5 & No \\
\midrule
LLaMA 3.2-1B$^*$ & Galaxy S24+ (NPU) & 45.2 & 85 & 1.2 & 3.8 & Good \\
LLaMA 3.2-3B$^*$ & Galaxy S24+ (NPU) & 22.5 & 165 & 2.4 & 6.2 & Good \\
Phi-4-mini & Galaxy S24+ (CPU) & 12.8 & 380 & 3.0 & 8.5 & Good \\
Qwen2.5-7B & Galaxy S24+ (CPU) & 5.4 & 920 & 5.5 & 18.2 & Marginal \\
\bottomrule
\end{tabular}%
}\\[4pt]
{\raggedright \small $^*$Reported with a 4-bit NPU backend (ExecuTorch/SpinQuant).\par}
\end{table}

The reported edge figures (Table~\ref{tab:edge_deployment}) draw clear
boundaries. An 8~GB Raspberry~Pi~5 is reported to run INT4 models up to about 3B
at marginally interactive speeds and sub-2B models comfortably. The Jetson Orin
Nano handles everything up to 9B at useful speeds. Snapdragon-class smartphones
with NPU delegation are reported to run models up to 7B, with interactive use
practical below 4B. Time-to-first-token, dominated by prefill, is the binding
constraint for conversational use: reported TTFT for a 3B INT4 model is around
4~s on the Pi but around 1~s on the Jetson. GQA and local--global attention
substantially reduce the reported KV-cache footprint, which otherwise dominates
memory at long context.

\subsection{Cost and Energy}
\label{subsec:reported_energy}

\begin{table}[ht]
\centering
\caption{Reported energy-per-inference estimates (joules per 1{,}000 completions, 512-token generation).}
\label{tab:energy}
\footnotesize\adjustbox{max width=\linewidth}{%
\begin{tabular}{lccccc}
\toprule
\textbf{Model} & \textbf{A100} & \textbf{RTX 4090} & \textbf{Laptop CPU} & \textbf{Jetson} & \textbf{RPi 5} \\
\midrule
Qwen2.5-0.5B & 180 & 320 & 410 & 280 & 612 \\
Qwen2.5-3B & 420 & 680 & 1,250 & 580 & 2,760 \\
Phi-4-mini & 510 & 780 & 1,480 & 650 & 3,180 \\
Qwen2.5-7B & 720 & 1,100 & 2,340 & 1,020 & --- \\
Phi-4 (14B) & 1,350 & 1,850 & --- & --- & --- \\
\bottomrule
\end{tabular}%
}
\end{table}

Reported energy figures (Table~\ref{tab:energy}) show that a datacenter GPU,
despite its high power draw, achieves the lowest energy \emph{per inference} for
any given model because of its throughput; an edge board consumes more energy per
inference but roughly 20$\times$ less absolute power, which favors it for
always-on, low-duty-cycle workloads. Normalizing accuracy by energy, a small
model on cheap hardware delivers markedly more benchmark points per joule than a
large model on a datacenter GPU, provided its accuracy clears the task threshold.
Table~\ref{tab:latency_compare} in Section~\ref{sec:deployment} translates these
into cost-per-token comparisons.

\subsection{Multilingual and Safety Reporting}
\label{subsec:reported_safety}

Reported multilingual scores (FLORES-200, XQuAD, MGSM) place Qwen2.5-7B ahead of
other SLMs, reflecting its large multilingual training corpus, with a wide gap to
English-centric models on non-English mathematical reasoning. Reported toxicity
evaluations (ToxiGen, RealToxicityPrompts) show substantially higher toxicity for
the smallest models (TinyLlama, sub-1B variants), consistent with the argument
that safety alignment is harder to preserve at reduced scale and is often absent
from research-grade small models. Both observations carry directly into
Section~\ref{sec:agentic}: an agent that routes non-English or safety-sensitive
calls to a small English-centric model inherits these gaps, which is an argument
for capability-aware rather than size-only routing.
